\chapter{Zwei weitere Eingaben: Mikrodynamik und deterministische Präparation}
\label{ch:newinputs13}
\begin{bild}{Aktualisierung v1.5: ein fester Puls ist nicht ein festes Detuning}
Die unten gezeigte Schranke eines einzigen verstimmten Pulses bleibt richtig. Eine endliche Folge positiver Pulse mit diskreten Q-abgeleiteten Belegungs-Z-Kicks und zwei Off-Intervallen erreicht dagegen vollständigen phasenrichtigen Transfer. Der fehlende Einzelpuls beweist also kein Verbot einer Kompositfolge. Siehe Kapitel~\ref{ch:execution15}.
\end{bild}
\begin{bild}{Aktualisierung v1.4}
Dieses Kapitel bewahrt die Ergebnisse und offenen Punkte von v1.3. Die bloße Belegungsbandtrennung ist inzwischen bis $|t|/\Delta=1/20$ bewiesen; kantenlokal gilt die neue Grenze $7\Delta/10$. Ein expliziter mikroskopischer Spektralfilter ersetzt im kontrollierten Sternvertrag die zuvor nur vorausgesetzte selektive Zieloperation. Ein eingeschränkter allordentlicher CAR-/Tensoradapter ist konstruiert. Details und unverändert offene native Herkunft: Kapitel~\ref{ch:audit14}--\ref{ch:fronts14}. Die hier erhaltene positive F4-Korrektur gehört zur zellintern geteilten Bank; sie darf nicht ohne Modellwechsel durch den kantenlokalen negativen Koeffizienten ersetzt werden.
\end{bild}
\section{Prüfurteil und Modellgrenze}
Die während der Konsolidierung zusätzlich eingereichten Texte N7 und N8 enthalten neue Resultate, nicht nur Zusammenfassungen. N7 berichtet die vollständige vierte Ordnung eines harten CAR-Referenzmodells, einen mikroskopischen Stern und einen Relaxationskanal. N8 konstruiert einen harten Tensorproduktkandidaten mit globaler Bandschranke und eine deterministische ideale Präparation.

Beide Statistikverträge werden getrennt. Die eigene Pfadrechnung verwendet harte lokale Tensorfaktoren, kohärente antisymmetrische Vertizes und harmonische Bosonen. Die Gleichheit mit einer bestimmten CAR-Phasenrealisierung oder dem nativen TFPT-Adapter wird nicht vorausgesetzt. Für den ausgeschriebenen Permutationsoperator wurde die neue große Singulettrechnung unabhängig reproduziert.

\section{Zwei richtige unitäre Vertizes, aber nicht derselbe Schritt}
N7 verwendet
\begin{equation}
U_W=\begin{pmatrix}P_+&-W^\dagger\\W&0\end{pmatrix}.
\end{equation}
Dieser Operator ist unitär, aber im Gegensatz zu $U_0$ aus Kapitel~\ref{ch:execution13} nicht involutiv:
\begin{equation}
U_W=U_0D_{\mathrm{occ}},\qquad
U_W^2=S\oplus(-I_6),\qquad U_0^2=I_{22}.
\end{equation}
Das ist kein Widerspruch, sondern eine verschiedene Phase des Rückwegs. Beim Recordaufbau ist der tatsächliche Adjungierte nötig:
$U_W^\dagger Q U_W$ reproduziert dieselbe Materie-Aufzeichnung wie $U_0QU_0$, unter Einschluss der Pointertensorfaktoren. Blindes Ersetzen beider $U_0$ durch $U_W$ wäre nicht derselbe phasentreue Prozess.

Allgemein gilt für endliche Operatoren mit gemeinsamem Eingangsraum
\begin{equation}
\widetilde K^\dagger\widetilde K=K^\dagger K
\quad\Longleftrightarrow\quad
\widetilde K=VK,
\end{equation}
wobei $V$ auf $\operatorname{ran}K$ isometrisch ist. Der Beweis definiert $V(Kx)=\widetilde Kx$: Gleiche Gram-Matrix macht dies wohldefiniert und skalarprodukttreu. Die Rückrichtung ist unmittelbar. Eine kohärente History muss daher eine isometrische Erweiterung des schon interferierenden Ausgangs sein, nicht eine kostenlose Markierung seiner einzelnen Eingangswege.

\section{Der vollständige vierte-Ordnungs-Operator}
Für $E_e=I-S_e$ lautet der eingereichte Referenzoperator
\begin{align}
H_{\mathrm{eff}}&=-\frac{t^2}{\Delta}\sum_eE_e
+\frac{t^4}{\Delta^3}F_4+O(t^6/\Delta^5),\\
F_4&=2\sum_eE_e+
\sum_{\substack{e<f\\e\cap f\ne\varnothing}}\{E_e,E_f\}
-2\sum_{\substack{e<f\\r(e)=r(f)}}E_eE_fT_{ef}.
\end{align}
$T_{ef}$ vertauscht die beiden vollständigen Paare. Die 40 Clebsch-Kanten ergeben 40 Einzelkanten-, 160 überlappende und 60 gleichlabelige disjunkte Beiträge. Jede Labelklasse ist ein Matching; gleiche Labels teilen keinen Endpunkt.

Die eigene Gegenrechnung rekonstruiert virtuelle Pfade in einer unnormalisierten Bosonbasis, in der Erzeugen Faktor eins und Vernichten Faktor $n$ trägt. Damit bleiben alle zurückkehrenden Nullboson-Matrixelemente rational. Die vierte Ordnung ist der gefaltete Beitrag $A^2$ minus die irreduzible Vierpfadantwort mit Energienenner $1\cdot2\cdot1$:
\begin{equation}
F_4=A^2-\tfrac12 B,\qquad
A=P_0VP_1VP_0,\quad B=P_0VP_1VP_2VP_1VP_0
\end{equation}
in Einheiten $t=\Delta=1$.

Diese Identität wurde auf allen 832 Eingangsspalten der kleinen Vergleichsmodelle geprüft: Dreierpfad, zwei disjunkte Kanten mit gleichem oder verschiedenem Vermittler und Viererstern. Zusätzlich wurden sechs vollständige C16-Matrixspalten einschließlich mehrfacher Farben berechnet. Das Weglassen der überlappenden Klasse wird durch eine Negativkontrolle erkannt.

Die Pfadklassifikation erklärt die Struktur: Auf einer Einzelkante gilt $E_e^2=2E_e$. Zwei überlappende Kanten können nicht gleichzeitig zwei Paare entfernen, daher bleibt ihr gefalteter Antikommutator. Disjunkte Kanten mit verschiedenem Label liefern nur den sich aufhebenden unabhängigen Beitrag. Geteilte Bosonlabels erlauben zusätzlich den Paartransfer $T_{ef}$. Im Clebsch-Matching können keine anderen gleichlabeligen Kanten dieselben Löcher anders paaren. Damit ist die Formel im benannten harten Tensorproduktvertrag begründet; eine native CAR-/Clock-Identifikation ist weiterhin ein eigener Satz.

\section{Die neue Singulettantwort ist unabhängig reproduziert}
Der gesamte 24024-dimensionale Singulett-Multiplizitätsraum wurde aus Youngs orthogonaler Darstellung aufgebaut. Ein neuer Lauf mit zwölf niedrigen Eigenvektoren ergab
\begin{align}
E_{0,s}/J&=11{,}045398337068436,\\
E_{1,s}/J&=11{,}561762122802579.
\end{align}
Beim zweiten Niveau wurden vier orthogonale Moden gefunden. Die größte Eigenvektorresiduen-Norm des Laufs liegt unter $7\cdot10^{-14}$. Das ist eine hochgenaue endliche Rechnung, kein algebraischer oberer Multiplizitätsbeweis.

Die vollständige Anwendung von $F_4$ liefert
\begin{align}
\langle0_s|F_4|0_s\rangle&=555{,}4885003638373,\\
\spec(P_{\mathrm{vier}}F_4P_{\mathrm{vier}})
&\approx\{583{,}292038666386^{[4]}\}.
\end{align}
Die Abweichungen von N7 betragen weniger als $8\cdot10^{-12}$. Mit $J=2t^2/\Delta$ und $\varepsilon=t/\Delta$ folgt
\begin{equation}
\frac{\Delta_s}{J}
=0{,}516363785734
+13{,}901769151274\,\varepsilon^2+O(\varepsilon^4).
\end{equation}
Bei $\varepsilon=1/20$ ist die erste korrigierte Näherung $0{,}551118208612J$, rund 6,73 Prozent über dem führenden Wert. Ein einzelner Ein-Prozent-Koeffizient erlaubt also keinen Schluss auf Ein-Prozent-Spektralfehler.

Nicht gerechnet wurden alle Nichtsinguletts, der vollständige mikroskopische Grundwert oder ein rigoroser Restterm höherer Ordnung. Die Referenznummern sind kein Beleg für einen bereits stabilen globalen Vielteilchengap.

\section{Eine globale Bandschranke, aber kein genauer Niedrigenergiefehler}
N8 definiert 16 harte Ortsfaktoren $\C\ket{\emptyset}\oplus\C^4$ und sechs Bosonmoden pro Vektorlabel. Bei Gesamtladung $\sum_sn_s+2\sum_{r,A}n_{r,A}=16$ ist der Sektor endlich. Für
\begin{equation}
H_0=\Delta\sum_{r,A}b^\dagger_{r,A}b_{r,A},\qquad
V=t\sum_{e,A}(b^\dagger_{r(e),A}K_{e,A}+K^\dagger_{e,A}b_{r(e),A})
\end{equation}
gilt $\|V_e\|\le4|t|$. Denn der antisymmetrische Umwandlungsblock $B_e=\sqrt2(b_1^\dagger,\ldots,b_6^\dagger)$ erfüllt $B_eB_e^\dagger=2n_rI$ und $n_r\le8$. Die selbstadjungierte Offdiagonalblocknorm ist dieselbe. Daher
\begin{equation}
\|V\|\le160|t|,\quad |t|/\Delta=1/640
\ \Longrightarrow\ 
\text{Bandabstand}\ge\Delta-2\|V\|\ge\Delta/2.
\end{equation}
Der führende Austausch ist dann $J/\Delta=1/204800$. Dies ist ein konservativer Existenzzeuge eines isolierten Vermittlerbands, nicht ein relativer Fehlerbound gegenüber $J$, ein interner Gap oder ein Vielzellenlimes. Die grobe Schranke zertifiziert $1/20$ nicht; sie widerlegt diesen Betriebswert auch nicht.

\section{Deterministische Präparation im idealen erweiterten Gattermodell}
Für $\chi=(\ket{01}-\ket{10})\otimes(\ket{23}-\ket{32})/2$ gilt $|\langle\Omega|\chi\rangle|^2=1/6$. Definiere $R_v(z)=I+(z-1)P_v$ mit $|z|=1$. Dann
\begin{equation}
[R_\chi(z)R_\Omega(z)]^2\chi=e^{i\gamma}\Omega,\qquad
\Re z=(3\sqrt5-7)/2.
\end{equation}
In der Basis $\{\Omega,\beta\}$ ist $\chi=s\Omega+c\beta$, $s^2=1/6$, $c^2=5/6$. Die unerwünschte Amplitude ist
\begin{equation}
\frac{c}{36}(z^4+14z^3+6z^2+14z+1)
=\frac{cz^2}{9}(u^2+7u+1),\quad u=\Re z.
\end{equation}
Sie verschwindet für die angegebene Wurzel exakt. Die ideale Phase ist $\phi=1{,}7172169856477322$. Es ist eine konkrete Anwendung bekannter \href{https://arxiv.org/abs/quant-ph/0005055}{Amplitudenverstärkung}, kein neuer allgemeiner Quantenalgorithmus.

$R_\Omega(z)$ kann unter kontrolliertem Zugriff auf den idealen Stern-Spektralfilter hergestellt werden: Ein Achtzweig-SELECT mit drei Kontrollqubits trennt $\Omega$ vom orthogonalen Raum am Kontrollausgang 000; eine Phase dort und vollständiges Rückrechnen setzt das Register zurück. Zwei $\Omega$-Phasen benötigen insgesamt zwölf kontrollierte Sternpotenzaufrufe, dazu zwei $\chi$-Phasen und deren Vorbereitung/Rückrechnung. Weitere Synthese- und Workspacekosten bleiben offen.

Ersetzt man jede der vier idealen selektiven Phasen durch einen unitären Operator mit Normfehler höchstens $\epsilon$, liefert Teleskopieren Zustandsfehler höchstens $4\epsilon$ und Zielwahrscheinlichkeit mindestens $1-16\epsilon^2$, bei exaktem Eingang. Kontroll- und Eingangsfehler müssen zusätzlich aufgenommen werden. Deterministisch im erweiterten Idealmodell heißt nicht kostenlos, autonom oder exakt in einem endlichen Clifford+$T$-Alphabet.

\section{Relaxation: konstruktiv möglich, aber nicht durch Bosonenverlust allein}
Gewöhnliche Vermittlerverluste $L_\mu=b_\mu$ lassen vollständig symmetrische Farbzustände mit Vermittlervakuum dunkel. Auf 16 Trägern hat $\mathrm{Sym}^{16}\C^4$ Dimension $\binom{19}{3}=969$. Dieser Raum verhindert einen eindeutigen globalen Attraktor dieser Verlustdynamik.

Für das ideale Vierträgersystem mit $G\ge0$, $\ker G=\C\Omega$, ist dagegen der Generator
\begin{equation}
\mathcal L_G(\rho)=-i(J/\hbar)[G,\rho]
+\kappa P_\Omega\Tr(G\rho)-\tfrac{\kappa}{2}\{G,\rho\}
\end{equation}
ein legitimer Lindbladgenerator. In einer Eigenbasis $G\ket j=g_j\ket j$ liefern $L_j=\sqrt{\kappa g_j}\ket\Omega\bra j$ die angegebene Form. Mit
\begin{equation}
B_t=e^{-(iJ/\hbar+\kappa/2)Gt}
\end{equation}
ist die vollständige Lösung
\begin{equation}
\rho(t)=B_t\rho(0)B_t^\dagger+
P_\Omega\Tr[(I-B_t^\dagger B_t)\rho(0)].
\end{equation}
Sie ist spurerhaltend und vollständig positiv. Bei Lücke $\delta$ gilt
$1-\Tr(P_\Omega\rho(t))\le e^{-\kappa\delta t}[1-\Tr(P_\Omega\rho(0))]$.
Die Sprünge enthalten den Zielzustand und eine ausgesuchte Reservoirkopplung. Deren Herkunft bleibt offen; der Satz beweist keine natürliche Kühlung durch die ursprünglichen Bosonenverluste.

\section{Der isolierte mikroskopische Stern: vollständiger kleiner Block}
Nach einem Paarabbau am Stern ist sein gemeinsamer Mittelpunkt leer; ein zweiter Paarabbau ist unmöglich. Für den geschlossenen Block aus 256 Materie- und 288 Vermittlerzuständen gilt
\begin{equation}
M^\dagger M=6I-2G_\star,\qquad
H_{\mathrm{mic}}=\begin{pmatrix}0&tM^\dagger\\tM&\Delta I_{288}\end{pmatrix}.
\end{equation}
Die untere effektive Matrix ist deshalb exakt
\begin{equation}
H_{\star,\mathrm{eff}}^{\mathrm{mic}}
=\tfrac12[\Delta I-\sqrt{\Delta^2I+4t^2(6I-2G_\star)}].
\end{equation}
Der volle 544-dimensionale Block wurde unabhängig gebaut und diagonalisiert. Bei $\Delta=1$, $t=1/20$ sind die niedrige Lücke $0{,}002433968751370$ und das nackte Materiegewicht des Grundzustands
\begin{equation}
w=\tfrac12(1+1/\sqrt{1+24(t/\Delta)^2})=0{,}9856429311786321.
\end{equation}
Die führende Lücke $J/2=0{,}0025$ ist nicht exakt. Ebenso ist der ideale Achtpunktfilter wegen der verschobenen Wurzelabstände kein automatisch exakter mikroskopischer Filter.

\section{Begonnenes Follow-up: die Präparation an das Mikromodell anpassen}
Für beliebige bekannte Zielüberlappung $a$ ist die Fehleramplitude nach zwei angepassten Schritten proportional zu
\begin{equation}
a^2(z-1)^4+3az(z-1)^2+z^2.
\end{equation}
Sie verschwindet bei
\begin{equation}
\cos\phi=1-\frac{3-\sqrt5}{4a},\qquad
a\ge\frac{3-\sqrt5}{8}.
\label{eq:matched-general}
\end{equation}
Die Polynomidentität wurde symbolisch geprüft. Für den \emph{angekleideten} Stern-Grundzustand und nackten Eingang $\chi$ gilt $a=w/6=0{,}16427382186310535$. Daraus folgt die angepasste Phase
\begin{equation}
\phi_{\mathrm{mic}}=1{,}7341107453033229.
\end{equation}
Verwendet man bei idealem Zugriff auf den angekleideten Zielprojektor dennoch die alte Phase, bleibt Infidelität $0{,}0001262823349012$. Die neue Phase beseitigt diesen Fehler im zweidimensionalen Idealvertrag. Der Zugang zum angekleideten Zielprojektor selbst wird dadurch \emph{nicht} konstruiert; genau diese noch fehlende selektive Operation ist der nächste Abnahmepunkt.

\section{Begonnenes Follow-up: warum ein einziger fester Puls nicht genügt}
Bereits der detunierte isolierte Paarblock hat maximale Transferwahrscheinlichkeit
\begin{equation}
p_{\max}=\frac{4g^2}{\Delta^2+4g^2}<1\quad(\Delta>0).
\end{equation}
Bei $g=\sqrt2t$, $t/\Delta=1/20$ sind das nur $1/51\approx1{,}96\%$. Eine einzelne freie Entwicklung dieses Blocks realisiert daher keinen vollständigen Materie-Vermittler-Tausch $U_0$.

Im ganzen Stern sind die Rabiabstände zu $M^\dagger M$-Eigenwert $a$ proportional zu $\sqrt{1+4\varepsilon^2a}$. Für $a=6,5$ und $\varepsilon=1/20$ ist ihr Verhältnis $\sqrt{106/105}$ irrational. Eine exakte Rückkehr des Vermittlers für \emph{alle} Materieeingänge nach einem einzigen gemeinsamen nichtverschwindenden Puls würde beide Sinusfaktoren gleichzeitig auf null setzen und damit ein rationales Frequenzverhältnis verlangen. Das ist ausgeschlossen.

Der Satz betrifft einen unveränderten freien Puls, nicht zusammengesetzte Pulse, zusätzliche Detuningskontrolle, selektive Messungen oder approximative Rückkehr. Er ist deshalb eine präzise Wegentscheidung: Die nächste Revision sollte kontrollierte Pulsfolgen oder einen nachweislich angepassten Spektralzugriff entwickeln, nicht denselben freien Puls nur länger laufen lassen.

\section{Ein kurzer positiver Kontrollkandidat für die nächste Revision}
Das Einzelpuls-Hindernis legt einen konkreten Ausweg nahe, sofern zusätzliche Kontrollen zugelassen werden. Im \emph{isolierten} resonanten Paar, also bei auf null gesetzter Verstimmung, sei
\begin{equation}
H_{\mathrm{res}}=g\begin{pmatrix}0&W^\dagger\\W&0\end{pmatrix},
\qquad g\tau/\hbar=\pi/2.
\end{equation}
Dann lautet seine exakte Entwicklung
\begin{equation}
U_{\mathrm{res}}=
\begin{pmatrix}P_+&-iW^\dagger\\-iW&0\end{pmatrix}.
\end{equation}
Mit dem relativen Viertelphasenpuls $D_{1/4}=I_{16}\oplus iI_6$ gilt
\begin{equation}
\boxed{D_{1/4}U_{\mathrm{res}}D_{1/4}=U_0.}
\end{equation}
Die volle 22-dimensionale Identität ist symbolisch geprüft. Zwei Belegungsphasen und ein resonanter Transfer genügen somit unter diesem Kontrollvertrag; es braucht keine unübersichtliche Gatterkaskade.

Das ist ein positiver bedingter Adapterkandidat, kein Nachweis nativer Verfügbarkeit. Man benötigt isolierte Kantenadressierung, veränderbare Verstimmung und einen relativen Belegungsphasenpuls. Die zentrale Clock liefert diesen Puls nicht automatisch. Während der Resonanz gilt außerdem nicht die zuvor benutzte große Vermittlerbandtrennung. Der nächste Test muss deshalb diese Kontrollen aus der Quelle begründen oder sie ausdrücklich als externe Steuerung mitführen und ihren Einfluss auf Zuschauer und Gesamtprozess berechnen.

\section{Arithmetik: bestätigt ist der Vertrag, nicht die offene Lösung}
Der singuläre Schur-Satz für $\begin{psmallmatrix}A&B\\B^\dagger&C\end{psmallmatrix}\ge0$ benötigt $A\ge0$, $\operatorname{ran}B\subseteq\operatorname{ran}A$ und $C-B^\dagger A^+B\ge0$. Das Gegenbeispiel $A=\operatorname{diag}(1,0)$, $B=(0,1)^T$, $C=2$ hat negatives Determinant trotz positivem Schurrest. Kompatible positive Blöcke ergeben eine positive dichte Fortsetzung nur für \emph{alle} Stufen, die tatsächliche Form und passende Stetigkeit. Die Identität der vollständigen signierten Weil-Form mit einer unabhängig definierten Norm bleibt offen. \href{https://arxiv.org/abs/2006.13771}{Connes--Consani} liefert den beschränkten archimedischen Hintergrund, keinen allgemeinen RH-Abschluss.

Für eine gerade unimodulare Gram-Matrix und $q(x)=x^TGx/2$ liefert die vollständige Fouriermessung des quadratischen Phasenzustands
\begin{equation}
\Pr(b)=(d/N)^8\,\mathbf1_{d\mid b_1,\ldots,b_8},\quad d=\gcd(t,N).
\end{equation}
Im Betragsquadrat setzt Charakterorthogonalität $tGh=0\bmod N$. Unimodularität reduziert dies auf $h=(N/d)z$; die verbleibende quadratische Phase ist ganzzahlig. Der Beweis gilt auch für gerade $N$. Die eigenen vollständigen Fourierverteilungen für $N=2,\ldots,6$ bestätigen ihn numerisch. Bei $d=1$ ist alles uniform; sonst ist nur der direkt aus $t,N$ berechenbare ggT sichtbar. Andere kohärente Anschlussoperationen sind damit nicht ausgeschlossen.

Faktorblinde Momentdurchläufe für $N=15,35,143,899,10403$ reproduzieren die Faktoren, aber benötigen jeweils $N$ ggT-Aufrufe. Der Umfang bleibt exponentiell in der Bitlänge. Die Katalogroute r647 bleibt entsprechend \texttt{STRUCTURAL\_MISMATCH}: Kodierung der Faktoren ersetzt keine effiziente Auslesung. Es wurde keine neue RH-Kampagne oder Faktorisierungsroute gestartet.

Die Aktualitätsprüfung des RH-Katalogs scheitert derzeit an einem fehlenden referenzierten Quellenordner; der konfigurierte Faktorgraph fehlt ebenfalls. Verfügbare Originale wurden gelesen. Eine vollständige aktuelle Kenntnis aller externen Versuche wird nicht behauptet.

\section{Die wichtigsten nächsten Fragen, einfach ausgedrückt}
\begin{enumerate}
\item \textbf{Kann die echte kleine Maschine die nötigen Handgriffe ausführen?}
Gesucht ist eine kontrollierte Kopplung, die Vermittler hin- und zurückführt und Aufzeichnung erhält. Ein fester freier Puls reicht im geprüften Stern nicht. Abnahme: beide Recordzweige und alle Eingänge, keine versteckten Zielzustände.
\item \textbf{Können wir den richtigen, angekleideten Zustand auswählen?}
Die passende Verstärkungsphase ist jetzt berechnet. Es fehlt der tatsächlich implementierte selektive Zielpuls einschließlich Genauigkeit und Kosten. Abnahme: Vorbereitung und Endmessung aus derselben Mikrodynamik.
\item \textbf{Bleibt die neue Spektralkorrektur bestehen, wenn alles mitgerechnet wird?}
Die ersten Singulettkorrekturen stimmen. Nächste Schritte sind höhere Restterme und Nichtsingulettkonkurrenz. Abnahme: kontrollierter Fehlerbereich, nicht bloß eine weitere Dezimalstelle.
\item \textbf{Wer liefert Reset, Kühlung und Anfangszustand?}
Ein maßgeschneiderter Relaxationskanal funktioniert mathematisch; gewöhnliche Verluste besitzen viele dunkle Zustände. Abnahme: konkrete Umgebung und Energie-/Entropiebilanz.
\item \textbf{Wie werden aus vielen Zellen Raum und Materie?}
Erst nach dem gemeinsamen kleinen Protokoll braucht es eine ausgewählte skalierende Familie, deren Korrelatoren Dimension, gemeinsame Ausbreitung und drei propagierende chirale Familien zeigen. Ein Index drei zählt ohne diesen Übertrag keine Familien.
\item \textbf{Wo entstehen Photonen und Gravitation tatsächlich?}
Ein isoliertes Vermittlerband ist kein masseloser Spin-1- oder Spin-2-Sektor. Abnahme: physikalische Pole, Eichidentitäten, zwei Gravitonhelizitäten und universelle Kopplung.
\end{enumerate}
Die ersten beiden Fragen wurden mit den neuen Phasen- und Pulsrechnungen begonnen. Keine dieser Fortschrittsmarkierungen schließt ein vollständiges T1--T8-Tor.
